\scanpage{42}
% archive sheet number: 41
\deleted{\diagram{\begin{tikzcd}[ampersand replacement=\&]
f^{-1}(U)\arrow[r,hook,"i"]\arrow[d,"f"]\&X\arrow[d,"f"]\\
U\arrow[r,"j"']\arrow[ur,"g"']\arrow[u,bend right=19,decorate,decoration={snake,amplitude=.4mm,segment length=2mm},"\deleted{\unsure{g}}"]\&Y
\end{tikzcd}}
\[
\begin{aligned}gf&=i\\fg&=j\end{aligned}
\]
$f$ morphisme birationnel \inserted{propre} de préschémas intègres $X,Y$

$U$ ouvert de $Y$ sur lequel $f^{-1}$ est défini [Si $Y$ est régulier en codim $1$, alors $U$ est le \unsure{complémentaire} d'\unsure{un|une} ens de codim $\geqslant2$], \deleted{\unsure{i.e. aussi les pts|i.e. ens des pts}} \unsure{i.e., aussi les pts|i.e., ensemble pts} ayant un voisinage $V$ tel que $f$ induise un isom $f^{-1}(U)\longrightarrow V$.

Soit $\xi\in\mathrm{Pic}(X)$ défini par un \inserted{\uline{cycle}} \uline{divisoriel} $D$ (\unsure{$X$ ouvert|$X$ connex|$X$ normal}), considérons \deleted{\unsure{un}} \deleted{le morphisme} $g^{-1}(\xi)\in\mathrm{Pic}(U)$, il correspond à un cycle divisoriel}
